\section{Model}
\label{sec.model}

We work in the static synchronous setting of~\cite{gkl24} and we keep some of
its notation. This section fixes the notation, states the assumptions, and
restates the facts about typical executions that the analysis uses.

\paragraph{Validators}

There are $n$ parties called validators. Up to $t$ may be Byzantine and
arbitrarily deviate from the protocol.

\paragraph{Rounds}

Time proceeds in rounds. Rounds are numbered $1,2,\dots$. For $r\le r'$, we
write $[r,r']$ for the set of rounds $\{r,\dots,r'\}$, and $[r,r']=\emptyset$
when $r'<r$.

\paragraph{Network}

We assume a synchronous network among the validators. A message sent by an
honest validator at round $r$ is received by every other validator at the
beginning of round $r+1$. The adversary can see the messages of honest parties
at any point of a round. At the end of each round it fixes the messages
delivered in the next round and their order. It may reorder messages, spoof
their source or add messages of its own, but it cannot drop a message.

\paragraph{Adversary}

We assume a \emph{Byzantine, rushing} adversary, which may corrupt up to $t$
parties. A rushing adversary sees every message sent by an honest validator
before it is delivered to any other party.

\paragraph{Random oracle}

We assume access to a random oracle $H \colon \{0,1\}^* \to \{0,1\}^\kappa$,
where $\kappa$ is the security parameter. Each validator is allowed $q$ queries
to the random oracle $H$ per round. The adversary is allowed $tq$ queries. A
query is successful if its output is below the target $T$. A single query is
successful with probability $p=T/2^\kappa$. A successful query made at round $r$
yields a block. We say that this block was \emph{computed at round $r$}.

\paragraph{Protocol}

Validators run the Bitcoin backbone protocol of~\cite{gkl24}. Every validator
keeps a local chain, initialized with the genesis block $\gen$. At round $r$ a
validator runs $\fn{maxvalid}$ on its local chain and on the chains it received
at the beginning of round $r$. Function $\fn{maxvalid}$ returns the longest
valid chain among its arguments and breaks ties towards the local chain. The
validator then makes its $q$ queries. On success it appends the new block to the
chain returned by $\fn{maxvalid}$. The resulting chain becomes its local chain
and is sent to all validators. Following~\cite{gkl24}, we say that a validator
\emph{has a chain at round $r$} if this was the local chain it started the round
with. The validator \emph{adopts a chain at round $r$} if this was the return
value of $\fn{maxvalid}$ at round $r$. The chain a validator adopts at round
$r$ is at least as long as the chain it has at round $r$. The chain it has at
round $r+1$ is the chain it adopted at round $r$, extended by at most one block
if the validator was successful at round $r$.

\paragraph{Clients}

Clients are parties that do not participate in the consensus protocol. A client
is only interested in confirming a ledger that is derived from the chain the
validators maintain. This ledger must be \emph{safe}. A client's connection to
the rest of the network is asynchronous; the adversary can delay messages for
any arbitrary, finite amount of time. Thus, the client cannot assume that it is
going to receive a message from an honest party within a certain number of
rounds. We say that a client \emph{receives a chain at round $r$} when the
adversary delivers that chain to it during round $r$.

Clients have access to a wall clock. They know the current round $r$, the
genesis block $\gen$, the number of validators $n$, the number of Byzantine
validators $t$, the probability that a round is successful $f$, the Chernoff
error parameter $\epsilon$ and the Chernoff concentration parameter $\lambda$.

\paragraph{Chains and the block tree}

A chain is a sequence of blocks, each of which contains the hash of its
predecessor. Every chain begins with the genesis block $\gen$. We write $\ch[i]$
for the $i$-th block of $\ch$, counting from $\ch[0]=\gen$, and we say that
$\ch[i]$ is at height $i$. We write $\ch[i{:}j]$ for the subsequence of $\ch$ at
heights $i,\dots,j-1$. We omit $i$ when the range starts at the beginning of
$\ch$, and $j$ when it extends to the end. Negative indices count from the end.
So $\head{\ch}$ is the last block of $\ch$, its \emph{tip}, and $\prune{\ch}{k}$
is $\ch$ without its $k$ last blocks. When $k\ge\len{\ch}$, $\prune{\ch}{k}=[]$.
For a real number $x$ we write $\prune{\ch}{x}$ for $\prune{\ch}{\lceil
x\rceil}$. We write $\len{\ch}$ to denote the length of $\ch$, which is the
number of blocks of $\ch$ including the genesis block. The height of $\ch$ is
the height of its tip, $\len{\ch}-1$. We write $\ch_1\preceq\ch_2$ when $\ch_1$
is a prefix of $\ch_2$.  That is, $\ch_1 \preceq \ch_2 \iff
\ch_2[{:}\len{\ch_1}] = \ch_1$. Every validator has a chain of length $1$ at
round $1$. A block is honest if it was computed by an honest validator and
adversarial otherwise. We regard the genesis block as an honest block computed
at round $0$.

\begin{definition}[Block tree]\label{def.tree}
    The block tree at round $r$, denoted $\bt_r$, is the set of all blocks
    computed in a round of $[0,r]$. Its edges are the hash references between
    these blocks. A block computed by the adversary belongs to $\bt_r$ whether or
    not it was sent. A chain $\ch$ is in $\bt_r$ if all of its blocks are in
    $\bt_r$.
\end{definition}

No copies occur in a typical execution, which we define below. So every block of
$\bt_r$ has a unique height and $\bt_r$ is a tree. Two facts about the block
tree follow from the protocol alone, and we use them throughout. The proofs are
in \Cref{appendix.model}.

\begin{lemma}[Membership]\label{lem.membership}
    Every chain held by a client at round $r$ is in $\bt_r$. Every chain
    adopted by an honest validator at round $r$ is in $\bt_{r-1}$.
\end{lemma}

\begin{lemma}[Fork choice]\label{lem.fork}
    Let $B$ be an honest block at height $i$, computed at round $r$. Then some
    honest validator has a chain of length greater than $i$ at round $r+1$.
    Every honest validator adopts a chain of length greater than $i$ at every
    round $r'>r$.
\end{lemma}

\begin{definition}[Common prefix of two chains]\label{def.meet}
    For chains $\ch_1$ and $\ch_2$, let $\ch_1\meet\ch_2$ be the longest chain
    that is a prefix of both.
\end{definition}

\paragraph{Random variables}

Fix a round $r$. Let $X_r=1$ if some honest validator computes a block at round
$r$, and $X_r=0$ otherwise. Let $Y_r=1$ if exactly one honest validator computes
a block at round $r$, and $Y_r=0$ otherwise. Let $Z_r$ be the number of
successful adversarial queries at round $r$. A round with $X_r=1$ is successful
and a round with $Y_r=1$ is a convergence opportunity. For a set of rounds $S$
we write $X(S)=\sum_{r\in S}X_r$ and define $Y(S)$ and $Z(S)$ in the same way.
We set $X_0=Y_0=Z_0=0$, so that a set of rounds containing the index $0$
contributes nothing at that index. The parameter $f=\mathbb{E}[X_r]$ is the
probability that a round is successful. From~\cite[Equation~(1)]{gkl24},
\begin{equation}\label{eq.rates}
    (1-f)\,pq(n-t) < f < pq(n-t) .
\end{equation}
Let $W_r=X_r-Y_r$, which is $1$ if two or more honest validators compute a block
at round $r$ and $0$ otherwise. Let $f_2=\mathbb{E}[W_r]$, and for a set of
rounds $S$ let $W(S)=\sum_{r\in S}W_r=X(S)-Y(S)$, with $W_0=0$. It is of order
$f^2$, so it is small compared with $f$. The proof is in \Cref{appendix.model}.

\begin{lemma}\label{lem.fsq}
    $f_2 < \dfrac{f^2}{2(1-f)^2}$.
\end{lemma}

\paragraph{Assumptions}

We write $\advr=t/(n-t)$ for the ratio of the adversarial to the honest query
budget, and we note that $1+\advr=n/(n-t)$.

\begin{assumption}[Honest majority]\label{asm.majority}
    There exist $\epsilon,\delta,f\in(0,1]$ such that $t\le(1-\delta)(n-t)$ and
    $3f+3\epsilon<\delta$.
\end{assumption}

We call $\delta$ the honest majority advantage. \Cref{asm.majority} is the
honest majority assumption of~\cite{gkl24}.

\paragraph{Typical executions}

Let $\lambda\ge2/f$ be an integer.

\begin{definition}[Typical execution]\label{def.typical}
    An execution is $(\epsilon,\lambda)$-typical if two conditions hold.
    First, every set $S$ of at least $\lambda$ consecutive rounds satisfies
    \[
        \begin{gathered}
            (1-\epsilon)f|S| < X(S) < (1+\epsilon)f|S|,
            \qquad Y(S) > (1-\epsilon)(1-f)f|S|,\\
            Z(S) < \left(\frac{\advr}{1-f}+\epsilon\right)f|S|,
            \qquad W(S) < \left(\frac{f_2}{f}+\epsilon\right)f|S| .
        \end{gathered}
    \]
    Second, no insertion, copy, prediction or guess occurs, in the sense
    of~\cite[Definition 4.3]{gkl24}.
\end{definition}

The upper bounds also hold for a set $[0,r]$ with $r\ge\lambda$, because
$X([0,r])=X([1,r])$, and likewise for $Y$, $Z$ and $W$, while the right-hand
sides increase with $|S|$. The definition is that of~\cite[Definition
4.4]{gkl24}, in the form of~\cite[Lemma 4.6]{gkl24}, with one bound added, the
bound on $W(S)$. Over any window of $\lambda$ or more rounds, every count stays
close to its expectation, within a margin set by $\epsilon$. Moreover, the hash
function behaves like a random function: nobody predicts a block, produces two
of the same, or inserts one into a chain it does not extend. Executions that
fail either condition are rare.

\begin{theorem}[Typicality]\label{thm.typical}
    An execution is $(\epsilon,\lambda)$-typical except with probability at
    most
    \(
        \typ=5L^2e^{-\Omega(\epsilon^2\lambda f)}+3Q^22^{-\kappa}+[(n-t)L]^22^{-\nu},
    \)
    where $L$ is the total run time of the system, $Q$ bounds the number of
    oracle queries, and $\nu$ is the min-entropy of the input contribution
    function.
\end{theorem}

For $\lambda=\Omega(\log^2\kappa)$ the probability $\typ$ is negligible in
$\kappa$. The theorem is~\cite[Theorem 4.5]{gkl24} with one Chernoff bound
added for $W(S)$, which is where its first coefficient is $5$ rather than $4$;
the proof is in \Cref{appendix.model}. All statements below are about typical
executions. The remaining facts of this section are restated
from~\cite{gkl24} without proof.

\begin{lemma}[{Pairing Lemma, \cite[Lemma 4.1]{gkl24}}]\label{lem.unique}
    Let $B$ be the block at height $i$ of a chain $\ch$, computed by an honest
    validator in a convergence opportunity. Then the block at height $i$ of
    any chain is either $B$ or adversarial.
\end{lemma}

\begin{lemma}[{Chain Growth Lemma, \cite[Lemma 4.2]{gkl24}}]\label{lem.growth}
    If at round $r$ an honest validator has a chain of length $m$, then at
    every round $r'\ge r$ every honest validator adopts a chain of length at
    least $m+X([r,r'-1])$.
\end{lemma}

\begin{lemma}[{Common Prefix, \cite[Theorem 4.10]{gkl24}}]\label{lem.commonprefix}
    Let $\ch_1$ and $\ch_2$ be chains adopted by honest validators at rounds
    $r_1$ and $r_2$ with $r_1\le r_2$. Then $\prune{\ch_1}{k}\preceq\ch_2$ for
    every $k\ge2\lambda f$.
\end{lemma}

\begin{lemma}[{\cite[Lemma 4.8]{gkl24}}]\label{lem.span}
    In a typical execution, any $k\ge2\lambda f$ consecutive blocks of a chain
    have been computed in more than $k/(2f)$ consecutive rounds.
\end{lemma}
